\chapter{The Flaws of the Similarity Assumption}\label{chap:flaws}
% ---- TAKEAWAYS ----
% - General-purpose representations are built for all of chemistry. On one dataset, most of their features barely vary.
% - Some representations give different molecules exactly the same vector, mostly stereoisomers. No model can tell those molecules apart.
% - A similarity or distance compares every feature, so features irrelevant to the task dilute it. Different representations also disagree about which molecules are close.
% - Methods that judge a molecule by its neighbors did worst under label noise: LNL methods that use a distance, and noise detection by neighbors.
% - What is needed is a compact set of features chosen for the task from several representations. Chapter 5 builds it: the hybrid.
% ---- END TAKEAWAYS ----
% ---- NOTES: The Flaws of the Similarity Assumption ----
% Takeaway: before the first section, say that one section reuses the confirmation report's LNL study.
% Paragraph 1: a short preamble, opening with \textit{Preamble.} as in Chapter~\ref{chap:robustness}.
% - Section \ref{sec:lnl_motivation} reuses the exploratory learning-with-noisy-labels study from the DPhil confirmation report.
% - The report is titled "Investigating the Effects of Noise on Molecular Property Prediction". It is unpublished (January 2026).
% - The text and numbers of that section are identical or similar to the report. The rest of the chapter is new.
% - Cite it once here as \citet{Punt2026confirmation}.
% - refs.bib entry needed: @unpublished{Punt2026confirmation, author = {Punt, Adelaide}, title = {Investigating the Effects of Noise on Molecular Property Prediction}, note = {DPhil confirmation report, University of Oxford}, year = {2026}}.
% ---- END NOTES ----

% CODE INPUT
This chapter asks whether the standard ways of comparing molecules suit the dataset and task at hand. General-purpose representations are built for all of chemistry, so we count how many of their features vary on one dataset. On hERG $K_i$, 1,815 of ECFP4's 2,048 features take one value in at least 95\% of the 1,412 molecules. Some representations also give different molecules the same vector, mostly stereoisomers, and no model can tell those molecules apart. A distance compares every feature, so features irrelevant to the task dilute it, and different representations disagree about which molecules are close. In an exploratory study of learning with noisy labels, all 20 configurations with the lowest retention used a distance, and noise detection by neighbors usually made the model worse.



\section{Introduction}
% [notes moved to thesis_notes_moved.md]

One of the fundamental assumptions in cheminformatics is that two structurally similar compounds display similar chemical properties \citep{Johnson1990}. While there are important exceptions, namely activity cliffs \citep{Stumpfe2012, vanTilborg2022}, this is a well-established concept. However, there is no single way to define similarity between compounds, and this is where the assumption begins to break down in practice. When defining similarity between two molecules, we typically rely on some pairwise mathematical function, or coefficient, that takes in two molecular representations as input. There's two varying factors here, the representation and the coefficient. \citet{Bajusz2015} found that among standard coefficient, the choice between them matters little. This question questions the underlying representation. 

General-purpose molecular representations are designed to work across chemical domains. Consequently, when considering data from a specific domain, many of their features show little to no variance. Some representations assign separate molecules identical feature vectors. Since common fingerprints do not encode stereochemistry \citep{orsi2024map4c}, they are often unable to differentiate stereoisomers, which can have substantially different measured activity \citep{Nguyen2006}. As these coefficients compare every feature, irrelevant features dilute the result. Different representations also disagree about which molecules are close based on the information they carry. 


The problem of whether neighboring molecules can predict chemical activity has been measure before. The MODelability Index (MODI) of \citet{Golbraikh2014} scores a binary dataset by the proportion of nearest-neighbor pairs that share an activity class, producing a single value for how strongly the similarity assumption holds on that particular combination of representation and data. It is not clear whether a low value is due to the data being intrinsically hard to model or to the representation not capturing adequate features for the task. Other work in this area relates to activity cliffs. \citet{Sheridan2020} compared experimental error, activity distribution and activity-cliff metrics as predictors of how well a dataset can be modeled, and found the cliff metrics the strongest of the three. However, they found that these metrics cannot separate real cliffs from  those produced by uncertain activity measurements. \citet{MedinaFranco2013} point out that some of these activity cliffs are artifacts of the chosen representation, similarity threshold, or noise in the data. Activity cliffs are distrinct from noise, as they carry some of the most valuable information about the interaction between structure and activity. None of these measures ask specifically about whether the specific features within these representations matter for the property at hand. 


There has been work that looks into how much the choice of representation changes similarity measure. \citet{Riniker2013} benchmarked 12 fingerprints for similarity across 3 public datasets and found high correlations between the, only two simple baseline fingerprints differed signficantly from the rest. The same was found for learned embeddings as neural embeddings gave negligible or no improvement over ECFP across 25 pretrained models and 25 datasets \citep{Praski2025}. \citet{Huber2026} did find a difference between the structure of the fingerprints. Measured using a graph-based structural reference, fingerprints that store raw counts improve agreement with the reference over folded fingerprints. The process of folding introduces bit collisions that distort similarity for molecules with many bit sets \citep{Huber2026}. Activity cliffs follow the same pattern, across 15 combinations of representation and similarity meausre on 3 datasets, the representation disagreed about which pairs form cliffs \citep{DabrowskiTumanski2026}. This supports the point \citet{MedinaFranco2013} make about cliffs as artifacts of their representation. Any method built on similarity like activity cliffs inherits this choice \citet{Wilson1972}. For the purposes of this thesis, these similarity measures are used across noise estimation and mitigation tasks. 

The purpose of this chapter is to establish the problem that Chapter~\ref{chap:method} sets out to solve. Section~\ref{sec:flaws:representations} asks for a given dataset, how much the features of standard representations vary and how many molecules map to identical feature vectors. 

Section~\ref{sec:flaws:similarity} looks at what happens when you used these features to quantify similarity, while Section~\ref{sec:flaws:consequences} further follows the consequences of using these similarity measures in downstream methods. This is supported by a sweep of LNL methods reported in Section~\ref{sec:lnl_motivation}.


\section{Flaws in standard base representations}\label{sec:flaws:representations} % CODE INPUT: label added, needs review
% [notes moved to thesis_notes_moved.md]

To quantify similarity between two compounds, we must first choose a representation to encode each compound. If the features within these representations are flawed, the similarity measure inherits them. Many standard representations are built to represent compounds across a wide array of chemical domains and be task-independent. However, chemical datasets often only cover a small region of the chemical space, leaving these general-purpose representations to describe chemistry which is absent from the dataset at hand. 

% TODO: come back to this


% CODE TODO [handled, see CODE INPUT blocks below]: Implement this and write it up. 
\subsection{Most features do not vary on a given dataset}\label{sec:flaws:constant} % CODE INPUT: label added, needs review
% [notes moved to thesis_notes_moved.md]
% ==== CODE INPUT: new, needs review ====
% Replaces the line: \todo[inline, color=yellow!40]{4.2.1: constant and near-constant feature counts (Table 4.1).}
% Numbers: scratchpad/thesis_code_todo/constant/constant_counts.json (scripts count_constant.py, mhg_dead.py, make_outputs.py in the same folder).
% Counts are on the saved feature matrices in rep_cache/<dataset>/<rep>.npy, all molecules of each dataset, no folds.
ECFP4 folded to 2,048 bits is built to describe any molecule, before any dataset or task is known \citep{rogers2010}. It is defined in Chapter~\ref{chap:background}. One dataset often comes from a few chemical series. On such a dataset most bits are off in almost every molecule. The features that matter for a task are a small subset of the ones that vary, and which subset depends on the task. This subsection counts only how many features vary at all. Whether the varying ones relate to the task is the question of Chapter~\ref{chap:method}.

We count two things per representation and dataset. A constant feature has one value in every molecule of the dataset. A near-constant feature has one value in at least 95\% of the molecules. The near-constant count includes the constant features. We count on all molecules of each dataset, using the feature matrices the full runs use. The datasets are hERG $K_i$, Caco-2, LogD and QM9, plus MDR1--MDCK efflux, human and rat liver microsome clearance, and the Lhasa hERG blocker training set. For each dataset we count its base representations, plus MACCS keys wherever they were computed.

On these datasets, the hashed substructure fingerprints have the most near-constant features. The learned embeddings have none, with one exception described below. ECFP4 on hERG $K_i$ (1,412 molecules) has 1,815 near-constant features out of 2,048. On the seven datasets where ECFP4 is a base representation, this count runs from 1,815 (hERG $K_i$) to 1,997 (QM9, 9,978 molecules). SMI-TED (768 features) and ChemBERT (256 features) on hERG $K_i$ have 0 near-constant features. GROVER, GraphMVP, MoLFormer, ChemBERTa-2, SELFormer and Mol2vec also have 0 on every dataset where they are base representations. The exception is MHG-GNN. Of its 1,024 features, 292 are zero in every molecule. They are the same 292 features on all seven datasets that use it, so they come from the saved MHG-GNN output, not from the data. We have not checked why the encoder gives zero there. The WL graph kernel also has 0 near-constant features. It keeps only labels present in at least 5\% of a dataset's molecules, so it is filtered on each dataset before it is counted. Every other count is in Tables~\ref{tab:constant_features} and~\ref{tab:constant_features_b}, and Figure~\ref{fig:constant_features} shows them as shares of each representation's features.

Which of the two counts matters depends on the distance. A bit that is off in every molecule does not change Tanimoto similarity, because bits off in both molecules are ignored. A bit that is on in every molecule is not harmless. It adds one to both the shared and the total bit counts, so it raises every Tanimoto similarity slightly. On hERG $K_i$ all 105 constant ECFP4 bits are off. On Caco-2, 1 of its 194 constant ECFP4 bits is on, and on LogD 1 of 39. On hERG $K_i$, 2 of the 24 constant MACCS features are on. After standardization, a constant feature adds zero to a Euclidean or Manhattan distance. Near-constant features do distort these distances. After standardization they take large values on the few molecules that have them. A bit that is on in 5\% of molecules becomes 4.36 on those molecules. A bit on in one of hERG $K_i$'s 1,412 molecules becomes 37.56 on that molecule. So the argument rests on the near-constant count, not the constant count. Section~\ref{sec:flaws:similarity} shows the effect on distances.

\citet[p.~112]{Dablander2024} states, without a count, that the least frequent substructures usually appear in only a few compounds. The same point recurs on p.~121. His statement is about substructures before hashing. Our count is of bits after folding to 2,048. Bit collisions (Chapter~\ref{chap:background}) also make a selected bit harder to read, because one bit can stand for several substructures. So a selected fingerprint bit is harder for a chemist to read than a named descriptor. This returns in Chapter~\ref{chap:method}, in the discussion of blueprints and what they tell a chemist.

If too few bits vary, different molecules can get identical vectors.

Tables~\ref{tab:constant_features} and~\ref{tab:constant_features_b} give, for each representation and dataset, the number of constant features and the number of near-constant features. The Features column gives the total number of features. The WL graph kernel's number of features differs between datasets, so its entries give the total as a third number.

\todo[inline]{TODO: every cell of this table must hold a count. The dashes are representations outside that dataset's base representations; handoff in HANDOFF\_fill\_constant\_features\_tables.md.}
\begin{table}[htbp]
\centering
\small
\setlength{\tabcolsep}{4pt}
\caption[Constant and near-constant features, part 1.]{Constant and near-constant features, part 1. Each entry is ``constant / near-constant'' features of that representation on that dataset. The number of molecules is under each dataset name. -- means the representation is not in that dataset's base representations. n/a means no saved MACCS features exist for QM9.}
\label{tab:constant_features}
\begin{adjustbox}{max width=\textwidth}
\begin{tabular}{lrcccc}
\toprule
Representation & Features & hERG $K_i$ & Caco-2 & LogD & QM9 \\
 & & (1,412) & (2,158) & (5,028) & (9,978) \\
\midrule
ECFP4 & 2,048 & 105 / 1,815 & 194 / 1,826 & 39 / 1,834 & 0 / 1,997 \\
Atom pairs & 2,048 & 201 / 977 & 563 / 1,239 & -- & -- \\
Topological torsions & 2,048 & -- & 1,179 / 1,822 & 1,059 / 1,818 & -- \\
Avalon & 2,048 & 50 / 985 & 98 / 1,128 & 30 / 1,076 & 4 / 1,380 \\
MACCS keys & 167 & 24 / 60 & 36 / 86 & 31 / 87 & n/a \\
WL graph kernel & 140--317 & 0 / 0 / 276 & 0 / 0 / 317 & -- & -- \\
PDV & 200 & 21 / 55 & 35 / 66 & -- & -- \\
Mordred & 1,613 & 242 / 336 & -- & 247 / 340 & 307 / 394 \\
GROVER & 4,800 & 0 / 0 & 0 / 0 & 0 / 0 & -- \\
MHG-GNN & 1,024 & 292 / 292 & 292 / 310 & 292 / 304 & -- \\
GraphMVP & 300 & -- & -- & 0 / 0 & 0 / 0 \\
MoLFormer & 768 & -- & 0 / 0 & 0 / 0 & 0 / 0 \\
SMI-TED & 768 & 0 / 0 & 0 / 0 & 0 / 0 & 0 / 0 \\
ChemBERT & 256 & 0 / 0 & -- & 0 / 0 & 0 / 0 \\
ChemBERTa-2 & 384 & -- & -- & -- & 0 / 0 \\
SELFormer & 768 & -- & -- & -- & 0 / 0 \\
Mol2vec & 300 & -- & -- & -- & 0 / 0 \\
\bottomrule
\end{tabular}
\end{adjustbox}
\end{table}

\todo[inline]{TODO: every cell of this table must hold a count. The dashes are representations outside that dataset's base representations; handoff in HANDOFF\_fill\_constant\_features\_tables.md.}
\begin{table}[htbp]
\centering
\small
\setlength{\tabcolsep}{4pt}
\caption[Constant and near-constant features, part 2.]{Constant and near-constant features, part 2. Entries read as in Table~\ref{tab:constant_features}. hERG blockers is the 7,772-molecule training set from Lhasa Limited.}
\label{tab:constant_features_b}
\begin{adjustbox}{max width=\textwidth}
\begin{tabular}{lrcccc}
\toprule
Representation & Features & MDR1--MDCK & HLM clearance & RLM clearance & hERG blockers \\
 & & (2,641) & (3,086) & (3,053) & (7,772) \\
\midrule
ECFP4 & 2,048 & 1 / 1,890 & -- & 1 / 1,887 & 0 / 1,840 \\
Atom pairs & 2,048 & 308 / 1,233 & 260 / 1,240 & 262 / 1,238 & 28 / 929 \\
Topological torsions & 2,048 & -- & -- & -- & 625 / 1,754 \\
Avalon & 2,048 & 5 / 1,032 & 8 / 1,044 & 8 / 1,045 & 0 / 866 \\
MACCS keys & 167 & 22 / 58 & 20 / 59 & 20 / 58 & 9 / 50 \\
WL graph kernel & 140--317 & 0 / 0 / 143 & 0 / 0 / 141 & 0 / 0 / 140 & 0 / 0 / 162 \\
PDV & 200 & 21 / 59 & 14 / 59 & 16 / 59 & -- \\
Mordred & 1,613 & 226 / 331 & 222 / 336 & 222 / 331 & 197 / 304 \\
GROVER & 4,800 & 0 / 0 & 0 / 0 & 0 / 0 & 0 / 0 \\
MHG-GNN & 1,024 & 292 / 292 & 292 / 293 & 292 / 293 & 292 / 292 \\
MoLFormer & 768 & -- & 0 / 0 & 0 / 0 & 0 / 0 \\
SMI-TED & 768 & 0 / 0 & 0 / 0 & 0 / 0 & -- \\
ChemBERT & 256 & -- & 0 / 0 & -- & -- \\
Mol2vec & 300 & -- & -- & -- & 0 / 0 \\
\bottomrule
\end{tabular}
\end{adjustbox}
\end{table}

MACCS keys have 167 features here. RDKit stores the 166 MACCS keys plus an unused bit 0, which is off in every molecule. That bit is one of the constant MACCS features counted above. On every dataset but QM9, all constant Mordred features are zero in every molecule. On QM9, one of the 307 is 1 in every molecule and the rest are zero. Missing or infinite descriptor values are set to 0 (Chapter~\ref{chap:method}, ``Representations by domain''). So some of these may be descriptors that failed to compute. We have not checked which.

% Uni-Mol v2 is left out of both tables. Its saved features are zero in every molecule on hERG Ki, LogD, MDR1-MDCK, HLM, RLM and the Lhasa hERG set (6 datasets, 768 features each). It is not in any dataset's base representations.
% The Lhasa hERG set appears once: rep_cache/herg and rep_cache/lhasa_herg hold the same 7,772 SMILES in the same order with the same labels.

\begin{figure}[htbp]
\centering
\includegraphics[width=\textwidth]{figures_code_input/constant_features.pdf}
\caption[Share of constant and near-constant features in each representation.]{Share of each representation's features that are constant (one value in every molecule), near-constant but not constant (one value in at least 95\% of molecules), and varying. One panel per dataset, one bar per representation. The number of features is in brackets after each name.}
\label{fig:constant_features}
\end{figure}
% ==== END CODE INPUT ====


\subsection{Molecules the representation cannot tell apart}
% [notes moved to thesis_notes_moved.md]

A standard fingerprint- or descriptor-based representation can produce the same feature vector for two separate molecules. In turn, models using this representation will produce a single prediction for the molecules with identical features. It should be noted that identical feature vectors within a dataset do not necessarily mean that they derive from the same underlying molecule. The dataset could just hold two different measurements of the same molecule twice, as is possible with data pooled across different sources.


Many fingerprints are computed without stereochemistry \citep{orsi2024map4c}. In this research, we build ECFP4 fingerprints with stereochemistry off while the 2D descriptor sets (PDV, Mordred) also ignore it such that for every dataset, each molecule they label a tie is a stereoisomer of another. Any representation read from the 2D graph merges stereoisomers, since the two share an identical graph. Separately, methods using the Weisfeiler--Lehman relabeling scheme are bounded by the 1-Weisfeiler--Lehman test explained in Chapter~\ref{chap:background}, so some non-isomorphic graphs also receive the same vector; the graph-kernel features tie 105 of 1,412 hERG $K_i$ molecules, 56 of which differ only in stereochemistry and 49 in more than that. Whether a learned embedding differentiates stereoisomers is up to its individual implementation. SMI-TED, a pretrained SMILES encoder, ties 52 hERG $K_i$, 633 LogD and 177 Caco-2 molecules, all stereoisomers. A tokenizer can drop stereochemistry before the model even sees the input string. The ChemBERTa-77M-MTR tokenizer used in Chapter~\ref{chap:robustness} gives 72 of 1,415 hERG $K_i$ molecules the same token sequence as another molecule. If the difference between stereoisomers was negligible, this wouldn't matter. However, stereoisomers can differ in measured activity \citep{Nguyen2006}, meaning these representations are all lacking potentially important chemical information.

% CODE TODO [handled, see CODE INPUT blocks below]: paragraphs 4/5/6 + table + figure. Basically finish it out
% ==== CODE INPUT: new, needs review ====
% Numbers: scratchpad/thesis_code_todo/ties/ties_results.json (script ties/ties.py; table rows ties/table.py; figure ties/figure.py, pair in ties/figure_pair.json).
% Paragraph 4: the measurement.
For each dataset and representation, we grouped the molecules whose stored feature vectors are exactly identical. We then gave every molecule in such a group one of the three reasons above. A molecule counts as a repeat if another molecule in its group has the same InChIKey. It counts as a featurization failure if its feature vector is all zeros. Otherwise it is a distinct molecule. We call a distinct molecule stereo-only when every molecule it shares a vector with has the same InChIKey connectivity block, so the two differ only in stereochemistry. No model is fitted.

Every molecule in a tied group must receive one prediction, and the prediction with the lowest squared error is the group's mean label. The lowest error any model can reach is therefore
\begin{equation}
\mathrm{RMSE}_{\min} = \sqrt{\frac{1}{n}\sum_{G}\sum_{i \in G}\left(y_i - \bar{y}_G\right)^2},
\end{equation}
where $n$ is the number of molecules, $G$ runs over the tied groups and $\bar{y}_G$ is the mean label of group $G$. A molecule in no tied group adds nothing to the sum. $\mathrm{RMSE}_{\min}$ is in the dataset's label units. The matching best $R^2$ is $1 - \mathrm{RMSE}_{\min}^2 / \mathrm{Var}(y)$. We also report the label spread inside each tied group, its highest label minus its lowest. For scale, repeat pK$_i$ measurements of one compound disagree by about $0.54$ log units \citep{Kramer2012}. A tied group whose pK$_i$ values differ by more than that is not explained by measurement noise alone. We counted on hERG $K_i$ (1,412 molecules after deduplication by InChIKey), LogD (5,028 molecules) and Caco-2 (2,158 molecules). Each dataset was counted on its final pool of representations and on its augmentation blocks.

% Paragraph 5: what was found.
No tie on any of the three datasets came from a repeated molecule or from a row of zeros. Every tie is a merge made by the representation. Table~\ref{tab:flaws:ties} lists every representation that ties at least one molecule. Atom pair, PDV, Mordred and SMI-TED tie only stereoisomers, on every dataset whose pool contains them. ECFP4, Avalon and topological torsion tie mostly stereoisomers, plus a few molecules that are not stereo-only. ECFP4 on hERG $K_i$ ties none of the latter. The most is topological torsion on LogD, with 44 such molecules. On hERG $K_i$, ECFP4, atom pair, PDV, Mordred and SMI-TED all split the 1,412 molecules into the same 26 groups, holding the same 56 molecules. So joining any of these five leaves all 56 tied. The graph kernel ties 49 hERG $K_i$ and 26 Caco-2 molecules that are not stereo-only. MACCS ties the most molecules on every dataset, and most of its tied molecules are not stereo-only. On hERG $K_i$ only 30 of its 377 tied molecules are. The 3D representations tie no molecules. Neither do the learned embeddings ChemBERT, GROVER, MHG-GNN, MoLFormer and GraphMVP, on any dataset whose pool contains them. % CHECK: the ChemBERT count comes from the cached chembert.npy (256 features). Whether that cache matches the current featurizer code is unverified.
This includes ChemBERT on hERG $K_i$ and on LogD. Joining every representation in a dataset's pool side by side leaves no molecule tied on any of the three datasets.

The lowest reachable error stays small next to the spread of the labels. On hERG $K_i$ it is at most $0.134$ pK$_i$ (MACCS), against a label standard deviation of $0.915$ pK$_i$. The spread inside one tied group can still be large. The largest ECFP4 group spread on hERG $K_i$ is $1.14$ pK$_i$, more than twice the repeat-measurement scale. Figure~\ref{fig:flaws:stereo_pair} shows that pair.

% Paragraph 6: consequence for activity cliffs and for neighbor-based methods.
The structure--activity landscape index of a pair of molecules $i$ and $j$ is
\begin{equation}
\mathrm{SALI}_{ij} = \frac{|y_i - y_j|}{1 - \mathrm{sim}(i,j)}
\end{equation}
\citep{Guha2008}. A tied pair has similarity 1, so its index is infinite whenever the two labels differ. A tie with different labels is therefore the most extreme activity cliff there can be. Some cliffs are made by the representation, not by the chemistry. A method that uses nearest neighbors sees a tied neighbor at distance 0 with a different label. % CHECK: sec:flaws:neighbors is not yet a label in thesis.tex; it comes from another CODE INPUT fragment.
It may then flag a correct label as noise (Section~\ref{sec:flaws:neighbors}).

% Table: one row per dataset x representation that ties any molecule.
Each row of Table~\ref{tab:flaws:ties} is one representation on one dataset. Only representations that give at least two molecules the same vector are listed.
\begin{table}[htbp]
\centering
\small
\setlength{\tabcolsep}{4pt}
\caption[Molecules that share a feature vector with another molecule.]{Molecules that share a feature vector with another molecule. \emph{Tied}: molecules whose vector is identical to at least one other molecule's (share of the dataset in brackets). \emph{Stereo-only}: tied molecules that share their vector only with their own stereoisomers. \emph{Best RMSE} and \emph{best $R^2$}: the lowest error any model can reach when every tied group gets one prediction, its mean label. \emph{Largest spread}: the highest minus the lowest label inside one tied group. RMSE and spread are in the dataset's label units, given after its name. (aug.) marks an augmentation block. Representations in a dataset's pool that are not listed tie no molecules.}
\label{tab:flaws:ties}
\begin{adjustbox}{max width=\textwidth}
\begin{tabular}{llrrrrrr}
\toprule
Dataset & Representation & Molecules & Tied & Stereo- & Best & Best & Largest \\
 & & & & only & RMSE & $R^2$ & spread \\
\midrule
hERG $K_i$ (pK$_i$) & ECFP4 & 1,412 & 56 (4.0\%) & 56 & 0.041 & 0.998 & 1.14 \\
 & Atom pair & 1,412 & 56 (4.0\%) & 56 & 0.041 & 0.998 & 1.14 \\
 & PDV & 1,412 & 56 (4.0\%) & 56 & 0.041 & 0.998 & 1.14 \\
 & Mordred & 1,412 & 56 (4.0\%) & 56 & 0.041 & 0.998 & 1.14 \\
 & SMI-TED & 1,412 & 56 (4.0\%) & 56 & 0.041 & 0.998 & 1.14 \\
 & Avalon & 1,412 & 66 (4.7\%) & 56 & 0.049 & 0.997 & 1.20 \\
 & Graph kernel & 1,412 & 105 (7.4\%) & 56 & 0.071 & 0.994 & 1.49 \\
 & MACCS (aug.) & 1,412 & 377 (26.7\%) & 30 & 0.134 & 0.978 & 1.80 \\
\midrule
LogD (log units) & ECFP4 & 5,028 & 641 (12.7\%) & 633 & 0.044 & 0.999 & 1.00 \\
 & Mordred & 5,028 & 633 (12.6\%) & 633 & 0.043 & 0.999 & 1.00 \\
 & SMI-TED & 5,028 & 633 (12.6\%) & 633 & 0.043 & 0.999 & 1.00 \\
 & Avalon & 5,028 & 647 (12.9\%) & 629 & 0.047 & 0.998 & 1.40 \\
 & Topological torsion & 5,028 & 675 (13.4\%) & 631 & 0.047 & 0.998 & 1.10 \\
 & MACCS (aug.) & 5,028 & 2,343 (46.6\%) & 437 & 0.220 & 0.965 & 2.80 \\
\midrule
Caco-2 (log units) & ECFP4 & 2,158 & 179 (8.3\%) & 177 & 0.027 & 0.996 & 0.74 \\
 & Atom pair & 2,158 & 177 (8.2\%) & 177 & 0.027 & 0.996 & 0.74 \\
 & PDV & 2,158 & 177 (8.2\%) & 177 & 0.027 & 0.996 & 0.74 \\
 & SMI-TED & 2,158 & 177 (8.2\%) & 177 & 0.027 & 0.996 & 0.74 \\
 & Avalon & 2,158 & 181 (8.4\%) & 177 & 0.027 & 0.996 & 0.74 \\
 & Topological torsion & 2,158 & 186 (8.6\%) & 177 & 0.029 & 0.996 & 0.74 \\
 & Graph kernel & 2,158 & 203 (9.4\%) & 177 & 0.028 & 0.996 & 0.74 \\
 & MACCS (aug.) & 2,158 & 804 (37.3\%) & 139 & 0.111 & 0.935 & 1.72 \\
\bottomrule
\end{tabular}
\end{adjustbox}
\end{table}

On hERG $K_i$, the mean label spread inside a tied group is $0.28$ pK$_i$ under ECFP4 and $0.39$ pK$_i$ under MACCS. Of the 26 ECFP4 groups, 24 hold molecules with different labels, and 6 differ by more than $0.54$ pK$_i$. Of the 153 MACCS groups, 138 hold different labels, and 38 differ by more than $0.54$ pK$_i$.

\begin{figure}[htbp]
\centering
\includegraphics[width=\textwidth]{figures_code_input/herg_ki_ecfp4_stereo_pair.pdf}
\caption[Two hERG $K_i$ molecules that differ at one stereocenter.]{Two hERG $K_i$ molecules that differ only at one stereocenter, the carbon bearing the amide nitrogen (R in A, S in B). Below each molecule are the ECFP4 bits it sets: the same 55 of 2,048 bits in both. The measured pK$_i$ values are $5.02$ and $6.16$. Any model on ECFP4 must give both the same prediction, and the one with the lowest squared error is their mean, $5.59$ pK$_i$. Each molecule is then $0.57$ pK$_i$ off.}
\label{fig:flaws:stereo_pair}
\end{figure}
% ==== END CODE INPUT ====


It is worth noting that some of these representations are more straightforward than others in terms of the underlying chemistry they reference. For example, the PDV carries very clear physicochemical properties such that if you examined one of it's features, you'd know exactly the type of chemical information it contains. Other representations like frozen embeddings, representing the last layer of a network, do not tell us what type of chemical information, if any, they relate to. 


\section{Flaws in similarity and other pairwise computations}\label{sec:flaws:similarity} % CODE INPUT: label added, needs review
% [notes moved to thesis_notes_moved.md]

We've already discussed some of the flaws of individual representations, but what happens when you compare two instanes of one of these representations? These pairwise comparisons, as described in Chapter~\ref{chap:background}, have the potential to violate the similarty assumption. When considering similarity and other pairwise comparisons, we discuss both the representation and the coefficient.  


The standard way of comparing two small molecules is the Tanimoto coefficient, which operates on any bit vector, applied to Morgan fingerprints \citep{LopezPerez2024}. Independently, the Tanimoto coefficient is described as the ``coefficient of choice'' for molecular similarity \citep{Bajusz2015}, while Morgan fingerprint, represented as ECFP4 in this research, is likewise one of the strongest fingerprints for the task, ranking among the strongest fingerprint for virtual screening and target prediction \citep{Capecchi2020}. ECFP4 and ECFP6 specifically, alongside topological torsion, rank among the best of 28 fingerprints tested for ordering diverse structures by similarity \citep{OBoyle2016}. However, it's not a one-size-fits-all solution. Across 96,456 human metabolites, 72.9\% had an ECFP4 nearest neighbor at Tanimoto distance 0 \citep{Capecchi2020}, again running into the issue of the representation lacking stereochemical information and more generally being unable to distinguish molecules. 

% CODE INPUT: references filled in (Equation~\ref{eq:tanimoto_kernel}; the labels sec:flaws:representations, sec:flaws:similarity, sec:flaws:consequences, sec:flaws:constant, sec:flaws:neighbors, sec:lnl_background were added at their headings). Needs review.
The Tanimoto coefficient is written out in Chapter~\ref{chap:background} as the ratio inside the Tanimoto kernel of Equation~\ref{eq:tanimoto_kernel}, specifically the number of bits set in both molecules divided by the number set in either. The corresponding distance is $1-T$ \citep{Bajusz2015}. Bits that are always off or always on across a dataset are counted in Section~\ref{sec:flaws:constant}. 


All a similarity coefficient can do is compare every feature, yet it cannot tell which features matter to the prediction task at hand. And the similarity assumption holds only relative to such a measure, which is fixed by the choice of representation. To estimate what this costs, we categorize a representation's features into those relevant to the task, $R$, and those that are irrelevant, $I$. We take the squared Euclidean distance as a sum over both:
\begin{equation}
d^2(\mathbf{x},\mathbf{y}) = \sum_{j \in R}(x_j-y_j)^2 + \sum_{j \in I}(x_j-y_j)^2 .
\label{eq:distance_split}
\end{equation}

Every irrelevant feature that varies adds one further term, so the irrelevant part grows linearly with $|I|$. After standardization each feature contributes roughly the same amount, so the share of the distance carrying signal is about $|R|/(|R|+|I|)$. When we use a general-purpose representation like ECFP4 on a dataset spanning a small portion of the chemiccal space, $|I|$ can be a lot larger than $|R|$. Features that remain constant across the dataset contribute nothing. Rather, the damage comes from features that vary but bear no relation to property in question. The worst case scenario involves a single rare bit set in only a handful of molecules. After standarization, uncommon bits take on larger values, outweighing common features that may be more relevant. This is why the near-constant count in Section~\ref{sec:flaws:constant} matters more than the constant count.


Larger representations, or higher dimensionality, compounds the problem. As the number of dimensions grow, the nearest and farthest molecules end up at nearly the same distance, meaning that these pairwise distances are less able to discriminate, as explained in Chapter~\ref{chap:background}. We can measure this effect by measuring the relative contrast, defined as the difference between the largest and smallest measured distance divided by the smallest, and hubness, a measure of how unevenely the molecules appear in other molecule's nearest neighbors \citep{Radovanovic2010}. These errors propogate along larger data pipelines, starting from the representation into similarity, from similarity into thes set of neighbors a molecule is compared against, and from there into a prediction of some kind. 

% CODE TODO [handled, see CODE INPUT blocks below]: paragraphs 6 + 7 + figure + table, finish this out


% ==== CODE INPUT: new, needs review ====
% Numbers: /Users/apunt/repos/KIRBy/scratchpad/thesis_code_todo/distance/results_all.json
% Scripts: compute.py, overlap.py, mcs.py, figure.py, collect.py in the same folder.
% Paragraph 5 result (cliff counts) is the first paragraph below; it can move up to Paragraph 5.
% HLM clearance (adme_hlm) is not in Table~\ref{tab:datasets}; drop its rows if it is not a thesis dataset.
% The cliff rule cited to Chapter 2 is so far only in that section's NOTES ("This thesis counts cliffs with a different rule"); its prose must state it.
We count activity cliffs with the rule of Chapter~\ref{chap:background}. A pair is similar at ECFP4 Tanimoto of at least 0.55. The activity gap is set per dataset so that 15\% of molecules are cliffs. The share moves in steps as the gap changes, so LogD reaches 14.5\%. The number of cliff molecules is therefore close to fixed by the rule, and the gap is what varies (Table~\ref{tab:flaws:distance_summary}). On hERG $K_i$, 211 of 1,412 molecules are cliffs at a gap of 1.76 p$K_i$. On QM9, only 1,093 pairs among 9,978 molecules reach a Tanimoto of 0.55. The gap that flags 15\% of QM9 molecules is then 0.05~eV, 0.06 times the median absolute deviation of the labels. Three datasets come from the Biogen set: rat liver microsome (RLM) clearance, MDR1--MDCK efflux and human liver microsome (HLM) clearance. On these the gap is 0.20 to 0.34 log units, 0.50 to 0.60 times the median absolute deviation. On QM9 and these three, the gap is smaller than the median absolute deviation of the labels. A gap of twice the median absolute deviation, the rule our code used before, flags 31 of 9,978 QM9 molecules and 3,529 of 5,028 LogD molecules.

The count also depends on the representation. We kept each dataset's gap and replaced ECFP4 with each other representation in the pool. For each, the similar pairs were its own closest pairs, as many as ECFP4 calls similar. On hERG $K_i$ the number of cliff molecules then ranges from 138 (atom pairs) to 267 (WL graph kernel). Only 57\% to 74\% of the pairs ECFP4 calls similar are among another representation's closest pairs. On QM9 that share falls to 1\% to 24\%. Every count is in Table~\ref{tab:flaws:distance_per_rep}.

Two different questions can be asked of a distance. The first is about all pairs: does distance track the difference in property across the whole dataset? Plotting descriptor distance against activity difference for every pair is the neighborhood-behavior plot of \citet{Patterson1996}. We summarize that plot with one number, the Spearman correlation between pair distance and absolute label difference. The second question is about nearest neighbors only: can each molecule's label be predicted from its five nearest neighbors? The prediction is the mean label of the five. The neighbors can carry signal even when the all-pairs correlation is near zero. In what follows, each claim says which of the two it is about.

We computed both for every representation in the pool of each of the seven regression datasets. Six are the small-molecule and QM9 datasets of Table~\ref{tab:dataset_reps}. The seventh is HLM clearance, which is not in that table. Representations with only 0/1 features were compared by Tanimoto distance, $1-T$. All others were compared by Manhattan distance on standardized features, leaving out features that are constant across the dataset. Manhattan distance is defined in the next subsection. For the all-pairs correlation, each dataset over 2,000 molecules was cut to a fixed random sample of 2,000 molecules (seed 0). That gives 1,999,000 pairs per dataset. hERG $K_i$ used all 996,166 pairs of its 1,412 molecules. For the neighbor prediction, all molecules were used, and each molecule's five neighbors came from other scaffold groups. No molecule is therefore predicted from its own scaffold.

Across the 70 combinations of dataset and representation, the all-pairs correlation ran from $-0.04$ (MACCS keys on MDR1--MDCK efflux) to 0.23 (PDV on the same dataset). It was below 0.10 in 44 of the 70. The five-neighbor prediction reached a Spearman correlation between 0.29 (Avalon on RLM clearance) and 0.89 (ChemBERTa-2 on QM9). In all 70 it was the larger of the two numbers. On hERG $K_i$, ECFP4 gives an all-pairs correlation of 0.06. Its five-neighbor prediction reaches 0.72, with an RMSE of 0.63 p$K_i$ against a label standard deviation of 0.92 p$K_i$. Table~\ref{tab:flaws:distance_summary} gives both numbers for three representations per dataset.

Figure~\ref{fig:flaws:herg_pairs} shows why the two numbers differ on hERG $K_i$. Of the 996,166 pairs, 93\% have an ECFP4 Tanimoto below 0.20, and 1.1\% reach 0.55. Across all pairs the median difference is 0.81 p$K_i$. Above a Tanimoto of about 0.4, the median falls to between 0.20 and 0.30 p$K_i$. Yet among the 774 pairs at 0.8 or above, 8.3\% still differ by more than 1 p$K_i$. Under Mordred, the closest 5\% of pairs have a median difference of 0.44 p$K_i$. Every other 5\% of pairs has a median between 0.78 and 1.11 p$K_i$. The signal sits in a thin tail of close pairs. The all-pairs correlation is set mostly by the other 95\% of pairs. A five-nearest-neighbor prediction looks only at each molecule's closest few.

A third view asks how often two representations agree on each molecule's five nearest neighbors. We compared ECFP4, one descriptor set and SMI-TED, a learned embedding that is in every pool. The descriptor set is Mordred, or PDV on Caco-2, where Mordred is not in the pool. HLM clearance's pool does not include ECFP4, so its ECFP4 was read from the same cache for this comparison only. The neighbor sets are the ones used for prediction. On hERG $K_i$ any two of the three share 0.47 to 0.52 of their five neighbors per molecule, on average. On the three Biogen datasets and QM9 they share 0.13 to 0.21, under one neighbor in five (Table~\ref{tab:flaws:neighbor_overlap}).

A fourth view asks whether fingerprint distance agrees with a structural distance. We sampled 500 random pairs per dataset (seed 0) and found their maximum common substructure with RDKit \citep{rdkit}. Ring atoms matched only ring atoms, only complete rings were kept, and each pair had a 2-second limit. The distance is $1 - n/(a+b-n)$, where $n$ is the number of atoms in the common substructure and $a$ and $b$ are the two molecules' heavy atoms. Pairs that hit the time limit were left out: 5 on hERG $K_i$, 8 on Caco-2 and 6 on LogD. The Spearman correlation between ECFP4 Tanimoto distance and this distance ran from 0.20 (hERG $K_i$ and HLM clearance) to 0.51 (Caco-2 and LogD).

A similarity can also be learned. \citet{Coupry2022} learn a molecular similarity from graph embeddings, but separately from any endpoint. It is still not chosen for the property being predicted. Earlier work did choose the similarity for the task. \citet{Zheng2000} predict each molecule's activity as the average of its $k$ most similar molecules, using Euclidean distance over descriptors. They select the subset of descriptors by simulated annealing to maximize the leave-one-out $q^2$. \citet{Itskowitz2005} extend this to optimize the descriptor subset, the number of neighbors and the neighbor weighting function together. Both select features for a nearest-neighbor model of one endpoint. The method in this thesis selects features across several representations. Every method in Section~\ref{sec:flaws:consequences} takes a similarity as input and inherits the dilution described here. ``Similarity in a hybrid'' (Chapter~\ref{chap:similarity}) tests whether a feature set chosen for the task gives a better similarity.

\begin{figure}[t]
\centering
\includegraphics[width=\textwidth]{figures_code_input/herg_ki_pair_distance_vs_dpki.pdf}
\caption[Absolute p$K_i$ difference against distance for hERG $K_i$ pairs.]{Absolute p$K_i$ difference against distance for all 996,166 pairs of the 1,412 hERG $K_i$ molecules. The line is the median in each bin and the band runs from the 25th to the 75th percentile. (a) ECFP4 Tanimoto similarity, in bins of 0.05; bins with fewer than 20 pairs are not drawn. The dashed line marks 0.55, the cliff rule's similarity. (b) Manhattan distance on standardized Mordred features, as a percentile of all pairs, in bins of 5\% of pairs. The left end holds the closest pairs, so the two panels run in opposite directions.}
\label{fig:flaws:herg_pairs}
\end{figure}

One row of Table~\ref{tab:flaws:distance_summary} is one dataset. The all-pairs columns give the Spearman correlation between pair distance and absolute label difference. The five-neighbor columns give the Spearman correlation between the five-neighbor prediction and the label.

\begin{table}[t]
\centering
\small
\setlength{\tabcolsep}{3.5pt}
\caption[All-pairs correlation, five-nearest-neighbor prediction and activity cliffs per dataset.]{All-pairs correlation and five-nearest-neighbor prediction for three representations, and activity cliffs, per dataset. Desc.\ is Mordred, except on Caco-2 efflux where it is PDV. Datasets over 2,000 molecules use a fixed sample of 2,000 molecules for the all-pairs columns. Neighbors come from other scaffold groups. Cliffs: ECFP4 Tanimoto $\geq 0.55$ and a label difference above the gap, which is set so that 15\% of molecules are cliffs. \textsuperscript{\textdagger}ECFP4 is not in this dataset's pool. Every representation is in Table~\ref{tab:flaws:distance_per_rep}.}
\label{tab:flaws:distance_summary}
\begin{adjustbox}{max width=\textwidth}
\begin{tabular}{lrrrrrrrrl}
\toprule
 & & \multicolumn{3}{c}{All pairs} & \multicolumn{3}{c}{Five neighbors} & \multicolumn{2}{c}{Cliffs} \\
\cmidrule(lr){3-5}\cmidrule(lr){6-8}\cmidrule(lr){9-10}
Dataset & Molecules & ECFP4 & Desc. & SMI-TED & ECFP4 & Desc. & SMI-TED & Molecules & Gap \\
\midrule
hERG $K_i$ & 1,412 & 0.06 & 0.07 & 0.09 & 0.72 & 0.72 & 0.69 & 211 & 1.76 p$K_i$ \\
LogD & 5,028 & 0.12 & 0.06 & 0.10 & 0.76 & 0.80 & 0.76 & 731 & 2.90 log \\
Caco-2 efflux & 2,158 & 0.08 & 0.18 & 0.14 & 0.59 & 0.62 & 0.57 & 324 & 1.21 log \\
RLM clearance & 3,053 & 0.02 & 0.07 & 0.04 & 0.36 & 0.49 & 0.46 & 458 & 0.34 log \\
MDR1--MDCK efflux & 2,641 & 0.02 & 0.16 & 0.10 & 0.51 & 0.58 & 0.59 & 397 & 0.20 log \\
HLM clearance & 3,086 & $-0.00$\textsuperscript{\textdagger} & 0.07 & 0.03 & 0.37\textsuperscript{\textdagger} & 0.49 & 0.45 & 462 & 0.27 log \\
QM9 & 9,978 & 0.19 & 0.18 & 0.13 & 0.88 & 0.84 & 0.85 & 1,495 & 0.05 eV \\
\bottomrule
\end{tabular}
\end{adjustbox}
\end{table}

One row of Table~\ref{tab:flaws:neighbor_overlap} is one dataset. The overlap columns give the fraction of a molecule's five nearest neighbors that two representations share, averaged over all molecules.

\begin{table}[t]
\centering
\small
\setlength{\tabcolsep}{4pt}
\caption[Agreement between representations per dataset.]{Agreement between representations, per dataset. Overlap: mean fraction of shared five-nearest-neighbor molecules, neighbors from other scaffold groups. MCS: Spearman correlation between ECFP4 Tanimoto distance and maximum-common-substructure distance on 500 random pairs (seed 0). Pairs that hit the 2-second limit are left out; the last two columns count the pairs used and left out.}
\label{tab:flaws:neighbor_overlap}
\begin{adjustbox}{max width=\textwidth}
\begin{tabular}{llrrrrrr}
\toprule
 & & \multicolumn{3}{c}{Five-neighbor overlap} & \multicolumn{3}{c}{ECFP4 against MCS} \\
\cmidrule(lr){3-5}\cmidrule(lr){6-8}
Dataset & Desc. & \makecell{ECFP4--\\Desc.} & \makecell{ECFP4--\\SMI-TED} & \makecell{Desc.--\\SMI-TED} & Spearman & Pairs & \makecell{Timed\\out} \\
\midrule
hERG $K_i$ & Mordred & 0.52 & 0.50 & 0.47 & 0.20 & 495 & 5 \\
LogD & Mordred & 0.38 & 0.41 & 0.41 & 0.51 & 494 & 6 \\
Caco-2 efflux & PDV & 0.38 & 0.41 & 0.38 & 0.51 & 492 & 8 \\
RLM clearance & Mordred & 0.15 & 0.21 & 0.18 & 0.28 & 500 & 0 \\
MDR1--MDCK efflux & Mordred & 0.15 & 0.20 & 0.17 & 0.26 & 500 & 0 \\
HLM clearance & Mordred & 0.15 & 0.21 & 0.18 & 0.20 & 500 & 0 \\
QM9 & Mordred & 0.13 & 0.18 & 0.21 & 0.46 & 500 & 0 \\
\bottomrule
\end{tabular}
\end{adjustbox}
\end{table}
% ==== END CODE INPUT ====


\subsubsection{Distance metrics for continuous molecular descriptors}
% [notes moved to thesis_notes_moved.md]

We now consider 8 commonly used distance metrics and how their formulation impacts down-the-line applications to distance-based noise detection. While binary fingerprints are compared with the Tanimoto coefficient, as described in Chapter~\ref{chap:background}, continuous features need a different measure. The performance of these 8 metrics in noise detection algorithms is compared in Section~\ref{sec:flaws:neighbors}.

\begin{itemize}
    \item Euclidean ($L_2$): standard geometric distance, $d(x, y) = \sqrt{\sum_i (x_i - y_i)^2}$
    \item Manhattan ($L_1$): sum of absolute differences, $d(x, y) = \sum_i |x_i - y_i|$, more robust to outlier features
    \item Chebyshev ($L_\infty$): maximum absolute difference, $d(x, y) = \max_i |x_i - y_i|$
    \item Cosine: angle-based, $d(x, y) = 1 - \frac{x \cdot y}{\|x\| \|y\|}$, invariant to magnitude
    \item Correlation: $d(x, y) = 1 - \rho(x, y)$, where $\rho$ is the Pearson correlation taken across the features of one pair of molecules, not across molecules
    \item Canberra: weighted $L_1$, $d(x, y) = \sum_i \frac{|x_i - y_i|}{|x_i| + |y_i|}$
    \item Bray--Curtis: ecological dissimilarity, $d(x, y) = \frac{\sum_i |x_i - y_i|}{\sum_i |x_i + y_i|}$ \citep{BrayCurtis1957}
    \item Mahalanobis: covariance-adjusted distance, $d(x,y)=\sqrt{(x-y)^\top S^{-1}(x-y)}$, where $S$ is the feature covariance estimated on the training molecules \citep{Mahalanobis1936}
\end{itemize}

All features are standardized to zero mean and unit variance prior to distance computation to ensure equal contribution across features with different relative scales. Bray--Curtis takes non-negative input, so for it the standardized matrix is shifted once to make its minimum zero. Canberra is computed on the standardized values directly. % TODO: REVIEW MODIFICATION -- formula checked against the code on 2026-09-29: tests/noise_mitigation_advanced.py:208-219 shifts Bray--Curtis only; Canberra only adds 1e-10. The Mahalanobis covariance matrix is singular here, as the features outnumber the molecules, so its pseudo-inverse is used.

Chapters~\ref{chap:method} and~\ref{chap:similarity} use four of these: Manhattan by default, Euclidean, cosine, and Minkowski with $p=0.5$, $d(x,y)=\left(\sum_i |x_i-y_i|^{0.5}\right)^{2}$, which breaks the triangle inequality and is therefore a dissimilarity rather than a metric. Feature weights from the hybrid are applied after standardization, as defined in Chapter~\ref{chap:similarity}.


% CODE INPUT: heading added. The notes block below had no \section, but the chapter introduction refers to it as Section~\ref{sec:flaws:consequences}. Needs review.
\section{Consequences for methods that rely on similarity}\label{sec:flaws:consequences}
% [notes moved to thesis_notes_moved.md]


We have identified flaws in the representations themselves and in the similarities computed over them. For a given molecule, its set of neighbors, built from one representation and one distance, is the input to several other processes such as predicting a molecule's activity from its nearest neighbors, detecting noisy labels by comparing a label against its neighbors' labels, defining an applicability domain by distance to the training set \citep{Tropsha2010}, grouping molecules into clusters for data splits \citep{Butina1999}, identifying activity cliffs, and, in federated learning, weighting each site's vote on a molecule by that molecule's distance to the site's own data \citep{Hanser2025}. If the distance is flawed, then every application of distance will suffer the consequences. When I was exploring LNL methods, I noticed that the ones that relied on the concept of molecular distance were unsuccessful, prompting me to look into the distance itself, as desribed in  Section~\ref{sec:lnl_motivation}.

Section~\ref{sec:lnl_motivation} reports the exploratory LNL study, in which every configuration that used a distance ranked among the worst. Section~\ref{sec:flaws:neighbors} turns to noise detection by neighbors, the most direct use of similarity in this thesis. 

% CODE TODO [handled, see CODE INPUT blocks below]: so I won't be just reporting my confirmation report. The reviewers won't see that. that material needs to be transferred over, and ideally shortened. this will take a lot... And I have enough there that you can go ahead and copy the langauge (maybe remove some) but refrain from adding your own flair unless it's absoltuely necessary. I've been ignoring every other instruction of yours to cite the confirmation report but now it's unavoidable. This needs fixing. And the output shouldn't be bullets it should just be a smaller version of what I wrote before with results. This should affect this section and the next.
\subsection{The exploratory learning-with-noisy-labels study}\label{sec:lnl_motivation}
% [notes moved to thesis_notes_moved.md]

% [OLD TEXT replaced by CODE INPUT below]
% Numerous studies have explored Learning with Noisy Labels (LNL) in the broader context of ML, as well as specifically for QSAR models. \citet{Song2022} provide a comprehensive review of Deep LNL, identifying five main approaches: (i) modifying architectures to better handle noise in training; (ii) employing regularization to reduce overfitting to noise; (iii) adjusting loss functions to be less sensitive to outliers or label flips; (iv) post-processing to adjust the weights for noisy data; and (v) sample selection to select presumably clean data to pre-train. \citet{Cortes2015} examined 12 different algorithms across 12 datasets and 10 noise levels, finding that no single algorithm was overall best at maintaining predictive performance in the presence of noise; certain algorithms perform better or worse at different levels of noise. This is a recurring theme in LNL research: there is no universal solution to handling noise with experimental endpoints.
% % TODO make this flow
% There are a wide variety of approaches to LNL, spanning sample selection, meta-learning, curriculum learning, and iterative refinement \citep{Song2022}. Cheminformatics data presents additional challenges in the form of high-dimensional molecular representations and dealing with a wide variety of chemical domains with various similarities and relationships. In this chapter, we explore a set of 12 LNL methods in the context of QSAR modeling. We examine the impact of uncertainty quantification, molecular distance metrics, and alternative loss functions in the hopes of improving performance under label noise. We tested 221 total configurations across LNL methods, loss functions, distance metrics, and molecular representations. Notably, we found that molecular representation choice explains even less variance under LNL methods. We also identified the following LNL techniques which showed promise on the tasks at hand: heteroscedastic GP, curriculum learning, and sample selection methods. The work was meant to connect the representation and uncertainty quantification insights found in Chapter~\ref{chap:robustness} with more classical noise mitigation techniques as a means of determining broader research directions. 
%
% One of the most prominent outcomes was that distance-augmented methods performed significantly worse than ones that did not incorporate any sense of distance or molecular similarity. An ANOVA test was run on the different adapted factors and shows that distance metric usage is the dominant factor in variance in $R^2$.. When ranking the performance of methods, all 20 worst-performing configurations with respect to noise robustness used distance metrics, not a single one achieved both high baseline ($R^2 > 0.6$) and high noise retention ($> 80\%$). This pattern was observed across standard distance types (Euclidean, Manhattan, Cosine, Mahalanobis) and a chemistry-aware metric (Tanimoto coefficient). 
%
% Building upon the exploratory study above, we decided to investigate alternative featurization techniques. There are multiple possible explanations behind the poor performance of distance metrics, first being that a singular molecular representation lacks sufficient information for distance calculations, with and possibly even without noise. These representations all encode different types of molecular information, PDVs encode global molecular properties, ECFPs contain substructure information, and MPNNs such as MHG-GNNs capture 3D structural features. These single sources of data may not be sufficient for robust similarity estimation. Additionally, some of these representations are sparse and contain a large number of features that are not relevant to molecules of a given domain. 
%
% The main issue we identified here is that when a large proportion of dimensions are uninformative, similarity measures become dominated by those uninformative features, causing exponential growth in the number of irrelevant pairwise interactions as dimensionality increases. This is what we sought to address by developing a representation with informative features and low dimensionality. We do so by creating a budget-constrained hybrid representation, which takes high-importance features from a variety of sources. These hybrid representations concatenate the top-ranked features by importance from a selection of base representations. 
% ==== CODE INPUT: new, needs review ====
% Transferred from the confirmation report (paper.tex, Chapter 4 "Exploring techniques for learning with noisy labels for QSAR", and the opening of Chapter 5), shortened, with the "Change:" bullets applied.
% TODO: Ch3 update -- re-check the Chapter 3 results quoted in this paragraph against the updated chapter.
This section reuses an exploratory study of learning with noisy labels (LNL) carried out earlier in this project. In Chapter~\ref{chap:robustness}, the choice of representation explained more of the variance in clean performance than the choice of model. The relationship flipped for performance lost under noise: there, the model had the stronger effect. We therefore asked whether training methods built for noisy labels improve noise robustness, and whether the representation still matters under them.

LNL methods, and the five groups of \citet{Song2022}, are described in Section~\ref{sec:lnl_background}. They have been studied across machine learning, including QSAR \citep{Cortes2015}. \citet{Cortes2015} examined 12 algorithms on 12 bioactivity datasets. They added Gaussian noise ranging from zero up to each dataset's range of pIC50 values. Gradient boosting tolerated noise worst. SVM and GP models with polynomial kernels were sensitive at high noise in some cases. The other algorithms degraded smoothly and linearly with the noise level. Most LNL methods were built for image classification with flipped labels \citep{Song2022}. The labels here are continuous, so these methods do not transfer automatically. Cheminformatics data adds further challenges: high-dimensional molecular representations and a wide variety of chemical domains.

The exploratory study tested 12 LNL methods on the HOMO-LUMO gap from QM9, using 10,000 molecules and scaffold splits. Gaussian noise with standard deviation $\sigma$ from 0 to 1.0, in steps of 0.1, was added to the training labels only. Test labels stayed clean. The methods were paired with five representations: ECFP4 (2,048 bits), Sort \& Slice \citep{Dablander2024}, PDV, canonical SMILES and randomized SMILES. Depending on the method, it could take up to four optional add-ons. These were the model's own uncertainty, a distance between molecules, an alternative GP kernel, and an alternative loss function. The distances were Tanimoto on ECFP4, Euclidean on PDV, and Mahalanobis. The GP kernels were Tanimoto, RBF and Mat\'{e}rn. In total we tested 221 configurations over method, loss function, distance and representation. The report grouped the 12 methods into five categories of its own: meta-learning, sample selection, curriculum learning, iterative refinement and GP extensions. These are not the five groups of \citet{Song2022}. Four of the 12 remove suspected molecules and then retrain: small-loss selection, confident learning, distance-based selection and multistage cleaning. These four are noise mitigation. The other eight are noise-robust training.

We measured noise robustness by retention: $R^2$ under noise divided by $R^2$ with clean labels, as a percentage. Methods that used a distance did worse than those that did not. All 20 configurations with the lowest retention used a distance. None of the configurations with a distance reached both a clean $R^2$ above 0.6 and a retention above 80\%. The failure sits in the distance step. The same methods without the distance add-on did not fail this way. Not all use of similarity failed, however. The heteroscedastic GP compares molecules through its kernel and still held up. The failure is specific to judging one label by its neighbors' labels.

The methods that held up estimate noise from the model's own loss or uncertainty. DivideMix \citep{li2020dividemix} kept the most, with 86.7\% retention (standard deviation 0.4\%). Small-loss selection \citep{Han2018} kept 84.2\% (standard deviation 2.1\%). The uncertainty curriculum \citep{Bengio2009} kept 82.7\% (standard deviation 3.5\%). The heteroscedastic GP \citep{liu2020} kept 79.0\%, averaged over 15 configurations (3 kernels and 5 representations). Its $R^2$ was 0.71 with clean labels and 0.56 at $\sigma = 0.6$, the noise level its retention was measured at. The report does not state the noise level for the other methods' retention. Regularization alone kept less: 73.0\% for Mixup and 74.0\% for sharpness-aware minimization. None of the methods that held up compares a label with its neighbors' labels. Under these methods, the choice of representation explained even less variance than in Chapter~\ref{chap:robustness}. Chapter~\ref{chap:similarity} compares twelve LNL methods on the current datasets.

The study used one dataset, one property and one noise type: Gaussian noise on every training label. Chapter~\ref{chap:similarity} repeats the question on more datasets and noise types. The distance step was also added to methods designed without one. A poor adaptation is therefore an alternative explanation. The recommendation of noise-robust training over noise mitigation (Chapter~\ref{chap:discussion}, Limitations) starts here.

There are two possible reasons for the poor performance of the distances. First, a single molecular representation may lack sufficient information for distance calculations, with and possibly even without noise. PDVs encode global molecular properties, ECFPs contain substructure information, and MHG-GNNs encode the molecular graph \citep{kishimoto2023}. These single sources of data may not be sufficient for robust similarity estimation. Second, some of these representations are sparse. They contain many features that are not relevant to the molecules of a given dataset. When a large proportion of features is uninformative, similarity measures become dominated by them. Each uninformative feature that varies adds one more term to the distance, so their share of it grows linearly with their number. The evidence here does not separate the two reasons. The second is the one we sought to address, with a representation of informative features and low dimensionality. The budget-constrained hybrid representation concatenates the top-ranked features by importance from a selection of base representations.
% ==== END CODE INPUT ====


\subsection{Noise detection by neighbors}\label{sec:flaws:neighbors} % CODE INPUT: label added, needs review
% [notes moved to thesis_notes_moved.md]

% [OLD TEXT replaced by CODE INPUT below]
% Overall, we see in agreement with conclusions of Section~\ref{sec:lnl_motivation} on the use of distance metrics in LNL methods, the concept of using a sample's neighbors to determine it's ground truth is flawed. It is likely all labels are too uniformly noisy for a methodology like this. Uncertainty-based methods show potential at lower noise, but their estimation performance degrades with higher noise. Sampling weighting methods generally proved effective, and one estimator using exponential sample weighting produced very strong recovery across both noise levels.
%
% Hybrid representations had mixed results. For sample removal, hybrids showed advantages at high noise, but struggled at low noise. They showed consistently positive recovery rates for neighbor-based estimation, which roughly falls in line with the idea that a low-dimension, high-quality information representation benefits neighbor-based reasoning. The feature budget in this experiment was restricted to 50-100 features, which may limit its advantages. This range was successful at standard molecular property prediction, but potentially not for certain noise mitigation or estimation tasks. Distance-based methods remain challenging under Gaussian noise regardless of representation, suggesting the core issue is the noise distribution rather than representation quality alone. This merits an investigation into mitigation of noise from other distributions or noise injection strategies.
% ==== CODE INPUT: new, needs review ====
% Transferred from the confirmation report (paper.tex, Chapter 5 "Noise Estimation" results, lines 1028-1146), shortened, with the "Change:" bullets applied.
Similar molecules should have similar labels. A label far from its neighbors' labels is therefore treated as probably wrong. The idea goes back to editing a nearest-neighbor training set by removing points whose label disagrees with their neighbors \citep{Wilson1972}. The wider literature on filtering noisy labels is about classification \citep{Brodley1999, frenay2014classification}. \citet{Martin2021} adapt fourteen classification noise filters to regression. They inject noise into the labels of 20 datasets and find that the filters improve the regression methods. Seven of the fourteen are nearest-neighbor filters: ENN, AENN, GE, BBNR, CNDC, CNN and RNN. For regression, each replaces the neighbors' majority class with the mean of the neighbors' labels \citep{Martin2021}. The other seven build ensembles of models and remove a sample by vote. Four of these (DF, EF, FMF and HRRF) include regression $k$-NN among their models \citep{Martin2021}. Their 20 datasets are general regression benchmarks, not QSAR.
% Source needed (from her NOTES, still open): a QSAR source for removing or down-weighting a molecule because its label disagrees with its nearest neighbors' labels.
This is the most direct use of the similarity assumption in this thesis. The scores, removal, weighting, smoothing and the recovery rate are defined in Chapter~\ref{chap:similarity}, ``Methods for learning with noisy labels''.

The method relies on two things at once. The distance must pick the right neighbors, and those neighbors' labels must be clean. The earlier sections of this chapter argue that the first fails. The second fails whenever noise is spread over the whole training set. A large gap between a label and its neighbors' labels can mean a noisy label, a noisy neighbor or a wrong set of neighbors. The score alone cannot tell which. Activity cliffs also look like noise to such a detector (Chapter~\ref{chap:background}, ``Molecular similarity and activity cliffs'').

The exploratory study also measured what happens when neighbors are used to clean labels. It used 5,000 QM9 molecules chosen at random, predicting the HOMO-LUMO gap, with a scaffold split. Gaussian noise was added to the training labels at $\sigma = 0.3$ and $\sigma = 0.6$. Each cleaning method was judged by retraining one random forest on ECFP4, on clean, noisy and cleaned labels. From these three fits we computed the recovery rate in $R^2$ (Chapter~\ref{chap:similarity}, ``Recovery rate''). Neighbor-based scores compared each molecule's label with the labels of its $k$ nearest training molecules. Euclidean, Manhattan and cosine distances were tried. The table below holds only the rows built on a single representation. Rows built on a hybrid used an allocation method that is no longer part of the thesis. This is a separate study from the one in ``Sample removal methods'' (Chapter~\ref{chap:similarity}). That one removed molecules, and used 10,000 QM9 molecules and eight distances. This one weighted and smoothed labels, and used 5,000 molecules and three distances.

In Table~\ref{tab:code_input_neighbor_recovery}, one row is one representation used to find neighbors. The neighbor-based scores were used as sample weights at strength $\alpha = 0.1$. With these weights, mean recovery was negative for every representation at both noise levels. Across all 860 cleaning configurations, 8 had positive recovery at both noise levels. Six of the 8 used a single representation, and 4 of those 6 were neighbor-based. Label smoothing toward neighbor or model estimates had positive recovery in 5.0\% of configurations at $\sigma = 0.3$ and 0.8\% at $\sigma = 0.6$. Mean recovery was negative at every weighting and smoothing strength tried, from 0.1 to 0.7.

\begin{table}[htbp]
\centering
\caption[Recovery in $R^2$ with neighbor-based sample weights, by representation.]{Recovery rate in $R^2$ when neighbor-based noise scores are used as sample weights ($\alpha = 0.1$), by the representation used to find neighbors. QM9 HOMO-LUMO gap, 5,000 molecules, random forest on ECFP4 as the downstream model. ``Positive'' is the share of configurations with recovery above zero.}
\label{tab:code_input_neighbor_recovery}
\begin{adjustbox}{max width=\textwidth}
\begin{tabular}{lcccc}
\toprule
 & \multicolumn{2}{c}{Mean recovery (\%)} & \multicolumn{2}{c}{Positive (\% of configurations)} \\
\cmidrule(lr){2-3}\cmidrule(lr){4-5}
Representation & $\sigma = 0.3$ & $\sigma = 0.6$ & $\sigma = 0.3$ & $\sigma = 0.6$ \\
\midrule
ECFP4   & $-3.0$ & $-7.3$ & 33.3 & 16.7 \\
PDV     & $-2.4$ & $-3.4$ & 50.0 & 27.8 \\
MHG-GNN & $-0.1$ & $-7.2$ & 38.9 & 5.6 \\
\bottomrule
\end{tabular}
\end{adjustbox}
\end{table}

Judging a label by its neighbors usually made the model worse, and rarely helped at both noise levels. It did not always fail. The best single configuration at $\sigma = 0.6$ was neighbor-based. It used PDV, Euclidean distance and the local variance ratio score, and recovered +17.2\%. Noise was added to every training label, however. A molecule's neighbors are then noisy too, so the neighbor mean is itself a noisy target. The report gives this as the reason label smoothing failed.

Overall, in agreement with Section~\ref{sec:lnl_motivation}, using a sample's neighbors to determine its ground truth is flawed here. It is likely that all labels are too uniformly noisy for a method like this. Uncertainty-based methods, which use no neighbors, show potential at lower noise, but their estimation degrades with higher noise. Sample weighting did better than smoothing but still usually hurt. It gave positive recovery in 14.2\% of configurations at $\sigma = 0.3$ and 4.7\% at $\sigma = 0.6$. Its mean recovery was $-22.9\%$ and $-20.9\%$. The configuration above, with exponential sample weighting, recovered +11.4\% at $\sigma = 0.3$ and +17.2\% at $\sigma = 0.6$.

The hybrid in that study was PDV plus MHG-GNN at 100 features, built by greedy forward selection. Greedy allocation is no longer part of the method. The current hybrid is built as described in Chapter~\ref{chap:method}, and its search sweeps budgets from 10 to 2,000 features. For neighbor-based scores, the old hybrid's mean recovery was $-0.9\%$ at $\sigma = 0.3$ and $-3.0\%$ at $\sigma = 0.6$. Its share of positive configurations was 44.4\% and 27.8\%. Distance-based methods remain challenging under Gaussian noise regardless of representation. This suggests the core issue is the noise distribution rather than representation quality alone. Chapter~\ref{chap:similarity} repeats the measurement with neighbors found in a hybrid, and under noise types other than Gaussian.
% ==== END CODE INPUT ====

% CODE TODO [handled, see CODE INPUT blocks below]: give me a sample to finish this out. Keep the relatively plain langauge style I've been working with, but I suspect this will change based on results a bit. Or not. I could just use some paragraphs on it. Don't go crazy, a conclusion is supposed to distill what's been done into something readable, something to remind the reader about the big picture before they walk away. Simple. 
\section{Conclusions: what a representation would need}
% [notes moved to thesis_notes_moved.md]
% ==== CODE INPUT: new, needs review ====
% Sample conclusion following NOTES Paragraphs 1-6 and the Table. No new evidence.
% Every number is copied from another CODE INPUT fragment of this chapter:
%   Paragraph 1: out/constant.tex (constant/constant_counts.json) and out/ties.tex (ties/ties_results.json).
%   Paragraph 2: out/distance.tex (distance/results_all.json).
% If those fragments change, these numbers must change with them.
This chapter asked three questions about the representations we use to compare molecules. Each was answered with a small measurement, not a new method.

The first question was how much of a standard representation varies on one dataset. Often very little does. On hERG $K_i$, 1,815 of ECFP4's 2,048 features take one value in at least 95\% of the 1,412 molecules (Section~\ref{sec:flaws:constant}). Several representations also give different molecules the same feature vector. ECFP4 gives 56 of the 1,412 hERG $K_i$ molecules a vector shared with another molecule, all of them stereoisomers (``Molecules the representation cannot tell apart''). No model can separate those molecules, whatever their labels.

The second question was whether distance follows the property. A distance sums over every feature. When only a few features matter for the property, the rest dominate the sum (Section~\ref{sec:flaws:similarity}). Across 70 combinations of dataset and representation, the Spearman correlation between the distance of a pair and its label difference was below 0.10 in 44. % CHECK: NOTES say keep the all-pairs vs nearest-neighbor contrast only if the full runs support it. The numbers below come from distance.tex, computed on the saved features, not from the full runs.
We also predicted each molecule's label as the mean of its five nearest neighbors' labels. That prediction's Spearman correlation with the label was higher in all 70. The useful information sits in a few close pairs, not in the distance as a whole.

The third question was what happens to methods that rely on distance. In the exploratory LNL study, the configurations that used a distance between molecules did worst under label noise (Section~\ref{sec:lnl_motivation}). Judging a label by its neighbors' labels usually made the model worse, and rarely helped at both noise levels (Section~\ref{sec:flaws:neighbors}). It did not always fail: one neighbor-based configuration gave the best recovery at the higher noise level. The noise was Gaussian and added to every training label, so each molecule's neighbors were noisy too. Chapter~\ref{chap:similarity} returns to other noise types.

There are two possible reasons the distance failed. One representation may carry too little information to judge similarity. Or a sparse representation may carry many features irrelevant to the dataset's molecules. Our evidence does not separate the two. The requirements below address the second.

A representation that avoids these flaws would need five things. Its features would be chosen for the task, using only the training labels of each fold. It would be compact: features that do not vary or do not matter are removed before any distance is taken or model is fit. The number of features kept would be searched, not fixed. It would draw from several representations, since each misses things the others capture. Plain concatenation meets this requirement but not the previous one (Chapter~\ref{chap:method}, ``Combining representations and the alternatives''). It would work both as a model input and as a space in which to compute similarity. And it would be traceable: each chosen feature leads back to its representation, and can be read chemically where the base feature can.

Feature selection cannot fix everything. Molecules that are identical in every base representation stay identical in the hybrid. A selected hashed bit still merges every substructure that folds into it. Selection also needs enough labeled molecules to estimate which features matter.

This is where the thesis began. Our hypothesis was that similarity as usually computed, ECFP4 with Tanimoto and its variants, carries many features irrelevant to the task. Methods built on it, such as noise detection and noise mitigation, suffer as a result. The first version of our method selected the task-relevant features of several representations, with the number of features set by the user. By a task we mean a supervised model fit on a training set. That version was built only to compute similarity. The selected feature sets then turned out to work as ordinary model inputs, without any similarity step. The selection itself became the method: the adaptive hybrid representation.

That leaves several questions for Chapter~\ref{chap:method}. How should a feature be scored (the importance method)? How should the budget be split across representations (the allocation method)? How should a candidate hybrid be scored (the internal evaluator)? How many features should be kept (the budget), and which downstream model should they feed? Features are chosen per representation, not from one pooled ranking, because some datasets prefer features from only the top few representations. Chapter~\ref{chap:method} makes that argument.

Table~\ref{tab:flaws:requirements} lists each flaw found in this chapter, what a representation would need to avoid it, and where the thesis handles it. One row is one flaw.

\begin{table}[htbp]
\centering
\small
\caption[Flaws of standard representations and similarity, and the requirement each sets.]{Flaws shown in this chapter, the requirement each sets, and where it is handled. The first column names sections of this chapter. The last names sections of Chapter~\ref{chap:method} unless another chapter is given.}
\label{tab:flaws:requirements}
\begin{adjustbox}{max width=\textwidth}
\begin{tabular}{P{0.34\textwidth}P{0.30\textwidth}P{0.28\textwidth}}
\toprule
Flaw (where shown) & Requirement & Where handled \\
\midrule
Features that do not vary (``Most features do not vary on a given dataset'') & Features chosen for the task, within a budget & ``Feature importance scoring'' \\
Molecules merged into one vector (``Molecules the representation cannot tell apart'') & Features drawn from several representations & ``Combining representations and the alternatives'' \\
Distance dominated by irrelevant features (``Flaws in similarity and other pairwise computations'') & Compact, with the number of features searched & ``The search space and why each axis is swept'' \\
Neighbor-based methods doing badly (``Consequences for methods that rely on similarity'') & Usable as a space for similarity & ``Similarity in a hybrid'' (Chapter~\ref{chap:similarity}) \\
\bottomrule
\end{tabular}
\end{adjustbox}
\end{table}
% ==== END CODE INPUT ====


 % Context for similarity and my motivating hypothesis: I think that molecular similarity as we often do it (ECFP4/Tanimoto or other variations of that) is flawed and applications that this concept of similarity (in my case noise detection and mitigation) suffer as a result
 %    - Take a look at ECFP4 as an example. Let's say you're working with molecules in a particular domain, how many of those features are going to useless? First, you can get numbers on in a given dataset how many features have no difference between the molecules. Some will differ, but will have nothing to do with the task at hand. There's definitely useful features in that set of 2048, but they're probably quite sparse, and when you take a similarity metrix that does a vector by vector comparison, and some of your features are irrelevant, you have exponentially more irrelevant pieces of information (is this true? phrase this better). I thought a much more efficient approach would be to select useful features for that task (I define task as some kind of supervised model fit on a training set) and you specify how many features you want, and you use that to calculate similarity. Now I implemented a proof-of-concept version of this and realized that these customized feature sets (hybrids) work as standard inputs to the model, irregardless of similarity. Then the main question was how on earth do I select the best features? How do I score how meaningful a feature is (feature importance), if I have lots of base reps, how do I know how many features from each to select (allocation strategies), once I have a candidate, how do I score it to see if its useful (internal evaluator). Then the other axes I test for the purposes of seeing if this works: does it work on different downstream evaluators, and what happens when I modify/control the budget (number of features).
 %        - Compare to other feature selection algorithms but explain why this is different - you could just pool all the features (and I've tried this, pull up my results on it), but it's not actually as effective as looking at them on a representation-level. It turns out if you have 10 reps you're considering, and you limit it to the top 2,3, 5, SOME datasets prefer that. Others prefer all 10. It's more complex than a simple feature selection algorithm if you really want to optimize it, and my results show why
 %        - This use case is specifically taking existing reps that say vastly different things about the molecule (PDV vs ECFP4 vs graph based embeddings vs sequence based embeddings) and tailor selecting a set of features from all of these to the very specific task at hand
 %    - Two concepts of using similarity: make a normal hybrid and use it for similarity OR modify the objective that the hybrid search optimizes for to make it similarity-aligned (called "similarity hybrids")
